\subsection{TV robustness, bit noise, and correlated thermal encodings}\label{ian:thermal-robust}
Status: proved unconditionally, as implications that use no Gibbs assumption unless one is stated: the perturbation estimates (Lemma~\ref{ian:thermal-robust-tv-estimates}); the TV-robust decoded compilation implication for Simon and period finding (Theorem~\ref{ian:thermal-robust-compilation}, Remark~\ref{ian:thermal-robust-decoder-scope}); the finite-energy bit-tube bounds, with arbitrary correlated auxiliary networks (Lemma~\ref{ian:thermal-bit-tube-obstruction}, Remark~\ref{ian:thermal-bit-noise-scope}); coherent block decoding and its reliability under the correlated-noise hypothesis \eqref{ian:thermal-error-domination} (Lemma~\ref{ian:thermal-coded-reliability}); all-spin hiding for independent noise at the level of a channel model (Proposition~\ref{ian:thermal-coded-hiding}) and the separation of reliability from hiding (Remark~\ref{ian:thermal-coded-channel-leakage}); quantum readout on a whole-table bit tube (Lemma~\ref{ian:thermal-global-bit-decoder}); the weight-only hardwired example (Proposition~\ref{ian:thermal-coded-promise-weight}): bounded degree-one fields around one public Simon table give decoded promise weight \(1-o(1)\) at every \(\beta\ge\beta_0\), but the decoded table equals that public table with probability \(1-o(1)\), so this region supplies no sector with \(\operatorname{Var}y\) bounded below and no hard sector for the Simon target; the correlated-interface obstruction (Theorem~\ref{ian:thermal-correlated-interface-obstruction}, Remark~\ref{ian:thermal-correlated-interface-scope}), which applies to the blockwise decoded table \(F_x=D_x(B_x)\) of codeword-symmetric word memories and, at inverse-polynomial sector weight, excludes \(r_m\le cm\) contacts per word for every fixed \(c<\log2/(\beta J)\); the verifier-leakage lemma (Lemma~\ref{ian:thermal-robust-verifier-leakage}). Proved conditionally: full-AA failure for Simon and period finding (Corollary~\ref{ian:thermal-robust-packages}), conditional on an actual physical sector of weight at least \(1/2\) whose decoded table and side information are within TV distance \(1/512\) of a product with the clean table law, exact coherent logical access, a joint whole-word simulation of every physical spin, auxiliary and reference spins included, and polynomial \(q,k\) with distinct physical dimensions. Open: a bounded-arity, bounded-degree thermal realization satisfying those hypotheses with unrestricted correlated auxiliaries, that is, a correlated network with sufficient word interfaces or without the stated word symmetries, no accessible target witness, the required random injective table law, and polynomial coherent logical access (Remark~\ref{ian:thermal-robust-frontier}). Literature: W.~Hoeffding, \emph{Probability inequalities for sums of bounded random variables}, J.~Amer.\ Statist.\ Assoc.\ \textbf{58} (1963), 13--30, Theorem~1 and the Chernoff optimisation in its proof, used in the proofs of Lemmas~\ref{ian:thermal-bit-tube-obstruction} and~\ref{ian:thermal-coded-reliability}.

\begin{definition}[Physical and clean laws]
\label{ian:thermal-robust-conventions}
All physical spins, including auxiliary and reference spins, belong to
\(\sigma\in\{-1,1\}^{N_m}\) and are queryable. The thermal class is
\[
 E_m(\sigma)=\sum_\alpha h_\alpha(\sigma_{S_\alpha}),
 \qquad |S_\alpha|\le k_0,\qquad
 \|h_\alpha\|_\infty\le J,\qquad
 \max_i|\{\alpha:i\in S_\alpha\}|\le d,
 \qquad
 \mu_{\beta,m}(\sigma)=Z_m^{-1}e^{-\beta E_m(\sigma)} ,
 \tag{38.1}\label{ian:thermal-robust-physical-bounds}
\]
where \(k_0,d,J,\beta>0\) are fixed and \(E_m\) is independent of
temperature. A sector law satisfies \(w_m\nu_m\le\mu_{\beta,m}\);
conditioning on an event of probability at least \(w_m\) is a special
case. We use \(\operatorname{Hard}\) and property \(\mathrm{(P)}\) from
Definition~\ref{ian:hard-sector-definition} and
Remark~\ref{ian:hard-sector-polynomial-quantifiers} in
\path{ian/13_hard_sector_definitions_forward.tex} and
\path{ian/15_hard_sector_sets_asymptotics.tex}.

The clean Simon law is Definition~\ref{ian:simon-law} in
\path{ian/30_simon_sector.tex}: \(M=2^m\), the nonzero mask is uniform,
and its \(M/2\) pairs receive uniformly random distinct \(m\)-bit
labels. Its target is \(Y=s_1\).

The clean period law is
Definition~\ref{ian:period-hsp-period-law} in
\path{ian/31_period_hsp_sectors.tex}: the period is uniform on the
prescribed even set of \(K\) primes, and the residue labels are a
uniform injection. Its target indicates the upper half of that
ordered prime set. In either case write \(\lambda_0\) for the clean
table law and \(y(F)\) for the target, extended by zero off its promise.
\end{definition}

\begin{lemma}[Two different perturbation estimates]
\label{ian:thermal-robust-tv-estimates}
If \(P,P'\) are laws at TV distance at most \(\xi\), and
\(U\in[0,1]\), \(V\in[-1,1]\) are fixed functions, then
\[
 \left|\operatorname{Var}_{P}U-\operatorname{Var}_{P'}U\right|
 \le\xi,\qquad
 \left|\operatorname{Cov}_{P}(U,V)
       -\operatorname{Cov}_{P'}(U,V)\right|
 \le5\xi.
 \tag{38.2}
\]
On one probability space, replacing \(U\) by \(U'\in[0,1]\) with
\(\mathbb E|U-U'|\le\rho\) changes its variance and its covariance
against a signed bounded test by at most \(2\rho\).
\end{lemma}

\begin{proof}
For the first assertion use
\(\operatorname{Var}U=\min_{c\in[0,1]}\mathbb E(U-c)^2\).
Every function being minimized takes values in \([0,1]\).
For covariance, the changes in
\(\mathbb E UV\), \(\mathbb E U\), and \(\mathbb E V\) are bounded
by \(2\xi,\xi,2\xi\), respectively. These are the estimates used in
Proposition~\ref{ian:hard-sector-soft-instance} of
\path{ian/17_hard_sector_instances.tex}.
The same-space estimates are
Lemma~\ref{ian:robust-compilation-perturbation} of
\path{ian/28_robust_compilation.tex}.
\end{proof}

\begin{theorem}[TV-robust decoded compilation]
\label{ian:thermal-robust-compilation}
Let \(w\nu\le\mu\) be an actual physical sector. Suppose that under
\(\nu\) there are a decoded table \(F\) and side information \(W\)
such that \(\sigma=\Psi(F,W)\) almost surely, with
\[
 \bigl\|\mathcal L_\nu(F,W)
       -\lambda_0\otimes\lambda_W\bigr\|_{\mathrm{TV}}\le\xi ,
 \tag{38.3}
\]
where \(\lambda_W\) is a law on the space of values of \(W\).
After \(W\) is fixed at any value in the support of \(\lambda_W\) or
of \(\mathcal L_\nu(W)\), each coordinate of \(\Psi(F,W)\) is computed
by a fixed, total algorithm making at most \(k\) adaptive whole-word
queries to \(F\). This includes every auxiliary and reference spin,
with the same \(W\) used for all answers.

Suppose also that a fixed coherent logical-coordinate oracle uses
at most \(q\) physical queries per call. It is defined on every
physical input and, \(\nu\)-almost surely, returns the corresponding
bit of \(F\) with its workspace restored.

Let a fixed clean-table quantum algorithm use \(C\) logical-coordinate
queries. Write \(q_0(F)\in[0,1]\) for its acceptance function on every
table. Assume that on the clean law
\(\mathbb E_{\lambda_0}|q_0-y|\le\delta_0\),
\(\operatorname{Var}_{\lambda_0}y=v_0\), and that every signed
whole-word tree of depth \(u\) has target covariance at most
\(\eta(u)\).

Then its physical implementation has a globally defined acceptance
function \(Q\), using at most \(T=qC\) physical queries, with
\[
 \begin{aligned}
 \mathbb E_\nu|Q-y(F)|&\le\delta_0+\xi,\\
 \operatorname{Var}_\nu Q&\ge
       a_\xi:=v_0-2\delta_0-\xi,\\
 \sup_{A\in\mathcal F_t}
       |\operatorname{Cov}_\nu(Q,A)|
       &\le\gamma_\xi:=
          \eta(kt)+2\delta_0+5\xi,\\
 \sup_{A\in\mathcal F_t}
       |\operatorname{Cov}_\nu(y(F),A)|
       &\le\eta(kt)+5\xi .
 \end{aligned}
 \tag{38.4}\label{ian:thermal-robust-compilation-bounds}
\]
In particular, the original hypothesis
\(F\perp W\), together with
\(\|\mathcal L_\nu(F)-\lambda_0\|_{\mathrm{TV}}\le\zeta\),
suffices with \(\xi=\zeta\) and \(\lambda_W=\mathcal L_\nu(W)\).
More generally, if the dependence distance
\(\|\mathcal L_\nu(F,W)
-\mathcal L_\nu(F)\otimes\mathcal L_\nu(W)\|_{\mathrm{TV}}\)
is at most \(\chi\), one may take \(\xi=\zeta+\chi\), again with
\(\lambda_W=\mathcal L_\nu(W)\).

Whenever \(0<\gamma_\xi<a_\xi\), these bounds give
\(\operatorname{Hard}(Q,\nu;t,a_\xi,\gamma_\xi)\).
The sufficient condition for \(\gamma_\xi\le a_\xi/4\) is
\[
 4\eta(kt)+10\delta_0+21\xi\le v_0 .
 \tag{38.5}\label{ian:thermal-robust-sufficient-condition}
\]
The sector weight remains \(w\).
\end{theorem}

\begin{proof}
This generalises
Theorem~\ref{ian:thermal-conditional-simon-compilation} in
\path{ian/36_thermal_realization.tex}, which assumes the exact clean
Simon law and side information independent of it.

For a physical tree \(A\), fixing \(W\) and composing its queries with
the coordinate simulators gives a whole-word tree
\(a(F,W)\in[-1,1]\) of depth at most \(kt\).
Under the comparison product law, conditioning on \(W\) preserves
the clean table law and its target mean. Consequently its target
covariance is at most \(\eta(kt)\). This is the side-information
reduction in the proof of
Theorem~\ref{ian:thermal-conditional-simon-compilation} and in
Lemma~\ref{ian:adaptive-character-side-information} of
\path{ian/29_adaptive_transcript_hiding.tex};
Theorem~\ref{ian:adaptive-transcript-duality} of the same file is used
only to pass from posterior-mean bounds to covariance bounds. Here the
required independence concerns the entire clean table.

On the actual sector, coherence and workspace restoration imply
\(Q(\sigma)=q_0(F)\). Under the comparison product law,
the target perturbation estimates give variance at least
\(v_0-2\delta_0\) and covariance at most
\(\eta(kt)+2\delta_0\).
Transport these fixed functions between the two laws using
Lemma~\ref{ian:thermal-robust-tv-estimates}.
The bounded function \(|q_0(F)-y(F)|\) gives the first estimate.

This comparison does not require the physical decoder to be correct
on inputs \(\Psi(F,W)\) generated only by the comparison law.
Its correctness is used on the actual sector; the comparison
functions are \(q_0(F)\) and the simulated classical transcript.
Finally, \(4\gamma_\xi\le a_\xi\) rearranges to
\eqref{ian:thermal-robust-sufficient-condition}.
\end{proof}

\begin{remark}[Decoder errors and the scope of TV approximation]
\label{ian:thermal-robust-decoder-scope}
A decoded marginal TV bound permits a \(\zeta\)-fraction of malformed
tables. The decoder is still one fixed circuit on the full physical
cube; it receives no promise test or sector oracle. Its average target
error becomes at most \(\delta_0+\zeta\) when side information is
independent. Hiding acquires an additive \(5\zeta\) term, so fixed
small TV error preserves constant hardness but does not establish
a vanishing hiding bound.

If the implemented acceptance function differs from \(q_0(F)\)
in sector \(L^1\) by at most \(\rho\), the variance lower bound in
\eqref{ian:thermal-robust-compilation-bounds} decreases by \(2\rho\),
the covariance upper bound increases by \(2\rho\), and the readout
error increases by \(\rho\). For example, if each logical-oracle
implementation has normalized diamond distance at most \(\kappa\)
from the required oracle, the hybrid argument gives
\(\rho\le C\kappa\).

Neither decoded TV closeness nor TV closeness of physical laws
establishes sector domination or its weight. Those remain separate
hypotheses.

Decoded marginal closeness alone also gives no all-spin hiding:
take an exact clean table and append the queryable auxiliary
\(y(F)\). Its decoded TV error is zero, while one physical query
has target covariance \(2v_0\).
\end{remark}

\begin{corollary}[Explicit full-AA-failure packages]
\label{ian:thermal-robust-packages}
Assume Theorem~\ref{ian:thermal-robust-compilation} with \(w\ge1/2\),
\(\xi\le1/512\), and exact coherent logical queries.

For Simon, let \(m\ge10\), \(M=2^m\), and set
\[
 T_{\mathrm S}=q\,m(m+9),\qquad
 t_{\mathrm S}=
       \left\lfloor\frac{\sqrt M}{64k}\right\rfloor ,
 \qquad
 T_{\mathrm P}=333qb,\qquad
 t_{\mathrm P}=
       \left\lfloor\frac{\sqrt K}{64k}\right\rfloor
 \tag{38.6}
\]
for the period law with \(K\ge1024\).
In either case the physical acceptance function satisfies
\[
 \operatorname{Hard}(Q,\nu;t,1/8,1/64),\qquad
 (w,a,\gamma)=(1/2,1/8,1/64),\qquad
 \ell_\mu(t,Q)\ge7/128 .
 \tag{38.7}\label{ian:thermal-robust-constant-package}
\]

For Simon, if \(q_m,k_m\) are bounded by fixed polynomials in \(m\)
and the physical dimensions are distinct along a subsequence, this
gives property \(\mathrm{(P)}\). The same conclusion holds for the
period family in Corollary~\ref{ian:period-hsp-full-aa}, with
\(q,k\) polynomial in \(\log R\).
\end{corollary}

\begin{proof}
For Simon use
Lemma~\ref{ian:simon-quantum-decoder}, and
Lemma~\ref{ian:simon-adaptive-posterior} together with
Theorem~\ref{ian:adaptive-transcript-duality} of
\path{ian/29_adaptive_transcript_hiding.tex}; the first two are in
\path{ian/30_simon_sector.tex}. They give
\(\delta_0=2^{-10}\),
\(v_0=1/4-1/[4(M-1)^2]\ge3/16\), and
\(\eta(u)=2\binom u2/(M-1-\binom u2)\).

For periods use
Lemma~\ref{ian:period-hsp-fourier-reconstruction} and
Lemma~\ref{ian:period-hsp-adaptive-hiding} in
\path{ian/31_period_hsp_sectors.tex}. They give
\(\delta_0=2^{-10}\), \(v_0=1/4\),
\(C=333b\), and
\(\eta(u)=u(u-1)/(K-u(u-1))\).

At the displayed depths, both hiding bounds are less than
\(1/2048\). Thus
\(a_\xi\ge3/16-1/512-1/512>1/8\), whereas
\(\gamma_\xi<
1/2048+1/512+5/512<1/64\).
The loss bound follows from
Proposition~\ref{ian:hard-sector-forward}.

With fixed \(r=16\) and \(K_{\mathrm P}=p_{\mathrm P}=q_{\mathrm P}=1\)
for the constants in property \(\mathrm{(P)}\), the inverse weight and
variance conditions hold. For every fixed budget constants \(A,B\),
\[
 \frac{t_m}{(T_m+16)^B}\longrightarrow\infty .
 \tag{38.8}
\]
For Simon this follows from exponential versus polynomial growth.
For periods it follows along the prime subsequence established in
Corollary~\ref{ian:period-hsp-full-aa}.
The stronger covariance bound in
\eqref{ian:thermal-robust-constant-package} implies the required
bound \(a/4\). These are conditional physical conclusions; the
compilation implication itself uses no Gibbs assumption.
\end{proof}

\begin{definition}[Encoded bit neighborhoods]
\label{ian:thermal-bit-neighborhoods}
Identify bits with physical signs. Let
\(C:\mathbb F_2^m\to\mathbb F_2^L\) be an injective linear code of
minimum distance at least \(\delta L\).
Choose a total decoder \(D\) correcting every error of weight
less than \(\delta L/2\).
It can be chosen equivariant:
\(D(z+C(a))=D(z)+a\), by fixing a minimum-weight representative
of each code coset.

For a physical word table \(z=(z_x)_{x\in\mathbb F_2^m}\), let
\(\mathcal S_m\) denote the clean Simon promise and define
\[
 \begin{aligned}
 A_\tau
   &=\{z:\exists f\in\mathcal S_m,\
       d_H(z_x,C(f(x)))\le\tau L\ \text{for every }x\},\\
 B_\tau
   &=\{z:\exists f\in\mathcal S_m,\
       \sum_xd_H(z_x,C(f(x)))\le\tau ML\}.
 \end{aligned}
 \tag{38.9}
\]
Thus \(A_\tau\subseteq\{D(z)\in\mathcal S_m\}\) when
\(\tau<\delta/2\), whereas \(B_\tau\) is a whole-table bit
Hamming tube.
\end{definition}

\begin{lemma}[Finite-energy bounds for bit tubes]
\label{ian:thermal-bit-tube-obstruction}
Under \eqref{ian:thermal-robust-physical-bounds}, put
\[
 \epsilon_*=\frac1{1+e^{2\beta dJ}},\qquad
 I(a\Vert p)=a\log\frac ap+(1-a)\log\frac{1-a}{1-p}.
 \tag{38.10}
\]
For \(0\le\tau<\epsilon_*\), and additionally
\(\tau<\delta/2\) for the first bound,
\[
 \begin{aligned}
 \mu_\beta(A_\tau)
   &\le(M-1)\exp\{-ML\,I(\tau\Vert\epsilon_*)/2\},\\
 \mu_\beta(B_\tau)
   &\le(M-1)\exp\{(m\log2)M/2
                         -ML\,I(\tau\Vert\epsilon_*)\}.
 \end{aligned}
 \tag{38.11}\label{ian:thermal-bit-tube-bounds}
\]
All auxiliary spins may form an arbitrary correlated network.
\end{lemma}

\begin{proof}
Lemma~\ref{ian:thermal-reconstruction-obstruction} in
\path{ian/36_thermal_realization.tex} gives conditional probability
at least \(\epsilon_*\) for either value of an unrevealed physical
bit, under conditioning on other coordinates.

Fix a mask and one representative of each pair. Condition on their
physical words and on all auxiliaries. On \(A_\tau\), each
representative decodes to its clean label, so it specifies the
required codeword at its partner. Reveal the \(ML/2\) partner bits.
Their mismatch indicators have conditional probabilities at least
\(\epsilon_*\). For \(h<0\), their exponential moment is at most
\((1-\epsilon_*+\epsilon_*e^h)^{ML/2}\).
The lower-tail exponential-moment bound gives
\(\exp\{-MLI(\tau\Vert\epsilon_*)/2\}\).
Sum over the \(M-1\) masks.

For \(B_\tau\), fix an entire clean encoded table and apply the same
argument to all \(ML\) bits. There are
\((M-1)(M)_{M/2}\le(M-1)M^{M/2}\) clean tables. A union bound gives
the second estimate.

This is a reconstruction of the conditional-probability argument.
The mismatch indicators need not be independent, so Theorem~1 of
W.~Hoeffding, \emph{Probability inequalities for sums of bounded
random variables}, J.~Amer.\ Statist.\ Assoc.\
\textbf{58} (1963), 13--30, is not applied directly. What is used is
the Chernoff optimisation over \(h\) in its proof, applied to the
iterated conditional exponential-moment bound above; it yields the
same relative-entropy exponent as that theorem.
\end{proof}

\begin{remark}[What finite energy does and does not specify]
\label{ian:thermal-bit-noise-scope}
The first bound in \eqref{ian:thermal-bit-tube-bounds} excludes
subthermal per-word tubes. The second excludes whole-table tubes
when \(LI(\tau\Vert\epsilon_*)>m\log2/2\) by a fixed margin
proportional to \(L\).

These bounds supply no exclusion for \(\tau>\epsilon_*\).
Finite energy supplies a lower conditional probability for either
physical bit value. It does not identify a latent clean table,
give an upper bound on corruption probabilities relative to that
table, or establish independence of errors. Conditioning on a
decoded table can itself involve the bit being examined.
Consequently a model in which bits are corrupted roughly
independently, each with probability about \(\epsilon_*\), requires
an additional, explicit channel hypothesis.
\end{remark}

\begin{lemma}[A sufficient correlated-noise hypothesis]
\label{ian:thermal-coded-reliability}
Suppose \(z_x=C(F(x))+E_x\). In a fixed enumeration of all error
bits, assume
\[
 \Pr(E_i=1\mid F,E_1,\ldots,E_{i-1})\le p<\tau<\delta/2 .
 \tag{38.12}\label{ian:thermal-error-domination}
\]
No independence among the errors is required. Then
\[
 \Pr(D(z)\ne F)\le
 \rho_m:=M e^{-LI(\tau\Vert p)} ,
 \qquad
 L\ge\frac{m\log2+a\log m}{I(\tau\Vert p)}
 \ \Longrightarrow\ \rho_m\le m^{-a}.
 \tag{38.13}\label{ian:thermal-coded-reliability-bound}
\]
A coherent whole-word decoder, with clean workspace, uses at most
\(2L\) physical queries and is defined on every physical input.
For a clean Simon latent table, the resulting Simon algorithm uses
\[
 T\le2L(m+9),\qquad
 \mathbb E|Q-y(F)|\le2^{-10}+\rho_m .
 \tag{38.14}\label{ian:thermal-coded-simon-readout}
\]
Also
\(\|\mathcal L(D(z))-\mathcal L(F)\|_{\mathrm{TV}}\le\rho_m\).
\end{lemma}

\begin{proof}
Iterating conditional exponential moments in
\eqref{ian:thermal-error-domination} gives, for each word and \(h>0\),
\(\mathbb E(e^{h|E_x|}\mid F)\le(1-p+pe^h)^L\).
The Chernoff optimisation over \(h\) in the proof of Theorem~1 of
Hoeffding's paper cited in the proof of
Lemma~\ref{ian:thermal-bit-tube-obstruction} gives
\(\Pr(|E_x|\ge\tau L\mid F)\le e^{-LI(\tau\Vert p)}\);
the error bits need not be independent, since only this
exponential-moment bound is used.
Union bound over addresses. On the complementary event every word
decodes correctly. The TV assertion follows from this coupling.

For coherent decoding, query the \(L\) bits of word \(x\), compute
the fixed total map \(D\) reversibly, XOR its output into the logical
answer register, undo that computation, and erase the queried word
with another \(L\) queries. This works with \(x\) in superposition.
The same construction returns an individual logical bit.
There is no query-cost assumption about the computation of \(D\).

On the simultaneous-correction event, \(m+9\) Fourier samples have
the clean Simon error at most \(2^{-10}\). Charge the remaining
event its full probability.
\end{proof}

\begin{proposition}[A hardwired public table: decoded promise weight without a hard target]
\label{ian:thermal-coded-promise-weight}
Fix \(J>0\) and constants \(0<\tau<\delta/2<1/4\), and fix one public
Simon table \(f_0\in\mathcal S_m\) for each \(m\). There are
\(\beta_0\), depending on \(J,\tau,\delta\), temperature-independent,
degree-one Hamiltonians of strength \(J\), and linear codes with
\(L=O(m)\) and relative distance at least \(\delta\), such that for
every \(\beta\ge\beta_0\)
\[
 \mu_\beta\{D(z)=f_0\}\ge1-\rho_m=1-o(1),
 \qquad\text{hence}\qquad
 \mu_\beta\{D(z)\in\mathcal S_m\}=1-o(1),
\]
with \(\rho_m\) from \eqref{ian:thermal-coded-reliability-bound}
exponentially small in \(m\).
The decoded table equals the one public table \(f_0\) with
probability \(1-o(1)\). Consequently every sector
\(w\nu\le\mu_\beta\) has
\(\operatorname{Var}_\nu y(D(z))\le\rho_m/w\), so this region supplies
no sector of inverse-polynomial weight with \(\operatorname{Var}y\)
bounded below, and no hard sector for the Simon target.
The proposition is a weight-only statement: an obstruction based on
decoded promise weight alone is impossible.
\end{proposition}

\begin{proof}
Put independent fields on the encoded bits of \(f_0\):
\[
 E_m(\sigma)=-J\sum_{x,j}
        (-1)^{C(f_0(x))_j}\sigma_{x,j},
 \qquad
 p_\beta=\Pr(\text{bit error})=\frac1{1+e^{2\beta J}} .
 \tag{38.15}
\]
Fix \(\beta_0\) with
\(p_{\beta_0}<\tau<\delta/2<1/4\).
Choose the code family once using these constants.
For every \(\beta\ge\beta_0\), the same Hamiltonian has independent
error probability at most \(p_{\beta_0}\).

For completeness, suitable linear codes exist by a direct random
matrix argument. For a uniformly random \(L\)-by-\(m\) binary
matrix \(G\), every \(Ga\), \(a\ne0\), is a uniform \(L\)-bit word.
Writing \(h(u)=-u\log u-(1-u)\log(1-u)\), a union bound gives
\[
 \Pr\{\exists a\ne0:\ |Ga|<\delta L\}
 \le 2^m e^{-L(\log2-h(\delta))}.
 \tag{38.16}
\]
Choose \(L=\lceil cm\rceil\), with
\(c(\log2-h(\delta))>\log2\) and
\(cI(\tau\Vert p_{\beta_0})>\log2\).
Such an injective code exists, and
\eqref{ian:thermal-coded-reliability-bound}, applied by
Lemma~\ref{ian:thermal-coded-reliability} with the deterministic
latent table \(f_0\) and \(p=p_{\beta_0}\), tends to zero
exponentially in \(m\). This gives
\(\mu_\beta\{D(z)\ne f_0\}\le\rho_m\).

The fields have strength \(J\) and degree one. They contain no
temperature parameter. Their decoded table concentrates on \(f_0\),
so this example does not approximate the prescribed random Simon
law and supplies no hard target distribution. Quantitatively, for
\(w\nu\le\mu_\beta\), the \(\{0,1\}\)-valued target satisfies
\(\operatorname{Var}_\nu y(D(z))
\le\nu\{y(D(z))\ne y(f_0)\}
\le\nu\{D(z)\ne f_0\}
\le\mu_\beta\{D(z)\ne f_0\}/w\le\rho_m/w\).
\end{proof}

\begin{proposition}[Independent noise preserves all-spin hiding]
\label{ian:thermal-coded-hiding}
Let \(F\) have the clean Simon law, and write
\(v_{\mathrm S}=\operatorname{Var}y(F)=1/4-1/[4(M-1)^2]\) and
\(\eta_{\mathrm S}(u)=2\binom u2/(M-1-\binom u2)\) for
\(\binom u2<M-1\), as in the proof of
Corollary~\ref{ian:thermal-robust-packages}. Suppose every physical
spin is generated from \((F,W)\), where \(W\perp F\), and is
simulatable using at most \(k\) whole-word queries after \(W\) is
fixed. For the encoded data bits with additive noise independent of
\(F\), one may take \(k=1\).

Let \(Q\) be the acceptance function of the Simon algorithm of
Lemma~\ref{ian:thermal-coded-reliability}, which uses the coherent
block decoder and satisfies \eqref{ian:thermal-coded-simon-readout}.
If the coherent block decoder succeeds simultaneously with
probability at least \(1-\rho_m\), then
\[
 \begin{aligned}
 \sup_{A\in\mathcal F_t}
      |\operatorname{Cov}(y(F),A)|
     &\le\eta_{\mathrm S}(kt),\\
 \operatorname{Var}Q
     &\ge v_{\mathrm S}-2(2^{-10}+\rho_m),\\
 \sup_{A\in\mathcal F_t}|\operatorname{Cov}(Q,A)|
     &\le\eta_{\mathrm S}(kt)+2(2^{-10}+\rho_m).
 \end{aligned}
 \tag{38.17}
\]
The analogous statements hold for periods with their clean hiding
bound and decoder.
\end{proposition}

\begin{proof}
Grant the entire error pattern and every independent seed to the
classical simulator. Fixing them leaves the clean table law intact.
Each encoded physical bit is a known function of one whole output
word. Compose adaptive queries exactly as in
Theorem~\ref{ian:thermal-robust-compilation}.
The remaining assertions follow from
Lemma~\ref{ian:thermal-coded-reliability} and
Theorem~\ref{ian:robust-compilation-theorem}.

For additional auxiliary spins, the required joint simulation is
an assumption about their actual law. It does not follow from
independence of the data-bit noise.
\end{proof}

\begin{remark}[Reliable correction need not preserve hiding]
\label{ian:thermal-coded-channel-leakage}
The upper-noise condition
\eqref{ian:thermal-error-domination} alone is insufficient for
hiding. For example, let all errors vanish except at one public
position of one word, where \(E_i=B\) and
\(\Pr(B=1\mid F)=p\,y(F)\). A code of distance at least three
corrects every such input. Reading that word reveals \(B\) through
its syndrome, and
\(\operatorname{Cov}(y(F),2B-1)=2p\operatorname{Var}(y(F))\).
This is constant for fixed \(p>0\), although correction is exact.
This channel example separates reliability from hiding; it is not
asserted to be a Gibbs realization.

For independent noise around a random clean table, reliability and
hiding therefore coexist at the level of a channel model. The
unproved thermal step is realizing that joint channel law, including
every auxiliary spin, by an allowed Hamiltonian.
\end{remark}

\begin{lemma}[Quantum readout on a whole-table bit tube]
\label{ian:thermal-global-bit-decoder}
If \(z\in B_\tau\) and \(\tau<\delta/8\), coherent block decoding
followed by the robust Simon word decoder recovers the mask with
error at most \(\varepsilon\), using
\[
 T\le
 2L\left\lceil
 \frac{32}{(1-8\tau/\delta)^2}
       \log\frac{2M}{\varepsilon}
 \right\rceil .
 \tag{38.18}
\]
The circuit is globally defined.
\end{lemma}

\begin{proof}
A wrongly decoded word requires at least \(\delta L/2\) bit errors
relative to the clean word. Thus at most
\(\theta M=(2\tau/\delta)M\) decoded words are wrong.
Here \(\theta<1/4\).
Apply Lemma~\ref{ian:thermal-simon-word-decoder} in
\path{ian/36_thermal_realization.tex}, replacing each raw whole-word
query by the coherent \(2L\)-query decoder.
Its exhaustive postprocessing is permitted by the query model.

Throughout this subsection, polynomial query complexity does not
assert polynomial gate complexity for an arbitrary code decoder or
for that exhaustive postprocessing.
\end{proof}

\begin{definition}[Word memories with a limited contact interface]
\label{ian:thermal-correlated-interface-class}
Partition the physical spins into word blocks
\(B_x=(z_x,u_x)\), \(x\in\mathbb F_2^m\), and an external auxiliary
network \(U\). Private auxiliary spins belong to their word blocks.
Suppose the Hamiltonian has the form
\[
 \begin{gathered}
 E=E_{\mathrm{net}}(U)+\sum_xE_x(B_x)
                       +\sum_\alpha h_\alpha,\\
 E_x(T_{x,a}B_x)=E_x(B_x),\qquad
 D_x(T_{x,a}B_x)=D_x(B_x)+a
       \quad(a\in\mathbb F_2^m),
 \end{gathered}
 \tag{38.19}
\]
where every \(T_{x,a}\) is a bijection of the word block.
At most \(r_m\) contact terms \(h_\alpha\) involve any given block.
All terms satisfy
\eqref{ian:thermal-robust-physical-bounds}.

The external network may be connected, frustrated, and arbitrarily
correlated. Contact terms may also couple different word blocks.
Any symmetry-breaking field inside a block counts among its contact
terms.

A natural instance is a linear-code memory whose internal energy
depends on parity-check syndromes, with codeword addition as
\(T_{x,a}\), coupled through a bounded number of contacts to a common
verifier network. Private gadget spins are allowed when their
bijection implements the stated energy symmetry.
\end{definition}

\begin{theorem}[A correlated-network obstruction]
\label{ian:thermal-correlated-interface-obstruction}
For the class in
Definition~\ref{ian:thermal-correlated-interface-class}, let
\(F_x=D_x(B_x)\), and put
\[
 a_m=e^{2\beta Jr_m},\qquad
 p_m=\frac{a_m}{a_m+M-1}.
 \qquad
 \mu_\beta\{F\in\mathcal S_m\}
       \le(M-1)p_m^{M/2}.
 \tag{38.20}\label{ian:thermal-correlated-interface-mass}
\]
In particular, if \(r_m=o(m)\), the right side is
\(\exp\{-(\log2/2-o(1))mM\}\).

If a sector \(w\nu\le\mu_\beta\) has decoded law within TV distance
\(\zeta<1\) of the clean Simon law, then
\[
 w(1-\zeta)\le(M-1)p_m^{M/2}.
 \tag{38.21}\label{ian:thermal-correlated-interface-sector}
\]
Consequently this class with \(r_m=o(m)\) cannot supply
Theorem~\ref{ian:thermal-robust-compilation} with polynomial
inverse sector weight in \(m\).
More quantitatively, if \(w\ge m^{-A}\) for fixed \(A\) and
\(1-\zeta\) is bounded below, a necessary condition is
\[
 2\beta Jr_m\ge2\log M-\log m-O(1).
 \tag{38.22}\label{ian:thermal-correlated-interface-necessary}
\]
\end{theorem}

\begin{proof}
Fix every physical spin outside \(B_x\). The map \(T_{x,a}\)
preserves its internal energy and changes at most \(r_m\) contact
terms, so its total energy change has absolute value at most
\(2Jr_m\).
It bijects the fiber \(D_x=b\) onto \(D_x=b+a\).
Hence the conditional partition sums of any two logical values
differ by a factor at most \(a_m\).
The largest logical conditional probability is therefore at most
\(a_m/(a_m+M-1)=p_m\).
This remains true after conditioning only on functions of the
other blocks, by conditional expectation.

Fix a mask and reveal one logical endpoint of every pair.
Successively requiring the other endpoints to equal their
representatives has conditional probability at most \(p_m\)
at each of \(M/2\) steps. Union bound over masks proves
\eqref{ian:thermal-correlated-interface-mass}.
Distinctness of pair labels only further restricts the event.
Every private and external auxiliary spin is included in these
conditional partition sums.

TV closeness implies
\(\nu(F\in\mathcal S_m)\ge1-\zeta\), giving
\eqref{ian:thermal-correlated-interface-sector}.
For \(r_m=o(m)\),
\(\log p_m=-m\log2+o(m)\).

Finally set
\(B_m=\log((M-1)/(w(1-\zeta)))=O(m)\).
Then \eqref{ian:thermal-correlated-interface-sector} implies
\(\log(1+(M-1)e^{-2\beta Jr_m})\le2B_m/M\).
Thus
\(e^{2\beta Jr_m}\ge(M-1)/(e^{2B_m/M}-1)
=\Omega(M^2/m)\), which proves
\eqref{ian:thermal-correlated-interface-necessary}.
\end{proof}

\begin{remark}[Scope of the construction attempt]
\label{ian:thermal-correlated-interface-scope}
The preceding theorem rules out codeword-symmetric word memories
with \(o(m)\) contacts per word even when their verifier auxiliaries
form an unrestricted correlated network. Increasing code distance
or block size does not remove this interface obstruction.
More precisely, when \(w\ge m^{-A}\) and \(1-\zeta\) is bounded
below, \eqref{ian:thermal-correlated-interface-necessary} reads
\(r_m\ge(\log2/(\beta J))m-O(\log m)\). It therefore also excludes,
for all large \(m\), linear interfaces \(r_m\le cm\) with any fixed
\(c<\log2/(\beta J)\).

The obstruction applies only to the blockwise decoded table
\(F_x=D_x(B_x)\) of
Definition~\ref{ian:thermal-correlated-interface-class}. It does not
rule out memories with \(r_m\ge(\log2/(\beta J))m-O(\log m)\)
contacts per word, encodings without the stated independent word
symmetries, or decoded tables that are not of this blockwise form.
The necessary bound \eqref{ian:thermal-correlated-interface-necessary}
is not a sufficiency statement. In particular, it neither constructs
a hidden pairing network nor proves that every sufficiently connected
network leaks its mask.

The pairwise and zero-field character realizations,
Theorems~\ref{res:pairwise-theorem} and
\ref{res:zero-field-theorem} in
\path{obstructions/local_codes/04_pairwise_all_spins.tex} and
\path{obstructions/local_codes/05_zero_field_pairwise.tex},
use their specific character symmetries.
Their all-spin reductions, including
Proposition~\ref{ian:hard-sector-local-instances} in
\path{ian/17_hard_sector_instances.tex}, do not establish the
Simon-table simulation hypothesis for a new correlated network.
Likewise, the soft-code Gibbs laws in
\path{obstructions/noisy_codes/03_fixed_hamiltonian_all_temperatures.tex}
do not supply a bounded-arity, bounded-degree realization.
\end{remark}

\begin{lemma}[Queryable verifier witnesses leak]
\label{ian:thermal-robust-verifier-leakage}
Let a physical depth-\(t_{\mathrm v}\) classical tree compute
\(\widehat Y\in[0,1]\), with
\(\mathbb E_\nu|\widehat Y-y(F)|\le e\).
Then the signed depth-\(t_{\mathrm v}\) test \(2\widehat Y-1\)
satisfies
\[
 \operatorname{Cov}_\nu(Q,2\widehat Y-1)
 \ge
 2\operatorname{Var}_\nu y(F)
       -2e-2\mathbb E_\nu|Q-y(F)|.
 \tag{38.23}
\]
Thus a verifier exposing the mask or period through a polynomial-query
classically decodable register fails the exponential-depth hiding
requirement, also under the small TV perturbations above.
\end{lemma}

\begin{proof}
This is Lemma~\ref{ian:thermal-exposed-witness} in
\path{ian/36_thermal_realization.tex}, applied to the decoded target.
Indeed,
\(\operatorname{Cov}(y,2\widehat Y-1)
=2\operatorname{Var}y+2\operatorname{Cov}(y,\widehat Y-y)\);
the last covariance has magnitude at most \(e\).
Replacing \(y\) by \(Q\) costs at most
\(2\mathbb E|Q-y|\).

For a readable mask register, the first mask bit gives the target.
Reading a period register determines its balanced target by free
classical computation. Redundant witness storage still leaks if
its classical readout uses polynomially many physical queries.
All such witness and decoding spins belong to the physical query
budget.
\end{proof}

\begin{remark}[Proved, conditional, and open conclusions]
\label{ian:thermal-robust-frontier}
The TV-robust implication and its explicit constants are proved.
Its thermal application remains conditional on actual sector
weight, coherent logical access, and a joint simulation of every
physical spin.

For bit noise, coherent decoding and hiding survive an independent
clean-table channel with the stated reliability bound. Polynomial
promise weight is achievable by bounded fields around a fixed
table, but there the decoded table equals that fixed table with
probability \(1-o(1)\), so the weight carries no hard target.
Those facts do not identify one Gibbs family satisfying
all three requirements. The low-radius bounds
\eqref{ian:thermal-bit-tube-bounds} and the positive example also
show why raw whole-word corruption obstructions cannot simply be
applied to thermally corrupted redundant words.

The correlated-interface class is refuted as a realization route
when \(r_m=o(m)\), and at inverse-polynomial sector weight whenever
\(r_m\le cm\) with fixed \(c<\log2/(\beta J)\). A local verifier with
a readable witness is separately excluded by leakage. A general
correlated network with sufficient word interfaces, no accessible
target witness, the required random injective table law, and
polynomial coherent logical access remains open for Simon and period
finding.
\end{remark}
